\chapter{Thin Lens Equation}

\section*{Introduction}

Your eye is a thin lens system. Light from the world passes through the cornea and lens, which bend it to form an image on the retina. The thin lens equation tells your brain how far away objects are based on how much the lens must bend the light. When you need glasses, it's because your eye's lens has the wrong focal length, and the glasses add a correction to bring the image into focus on the retina.

\section*{Fundamental Equation Used}

\[
\frac{1}{f} = \frac{1}{d_o} + \frac{1}{d_i}
\]
\section*{Assumptions}

\textbf{Assumptions:} The lens is thin (thickness $\ll$ focal length). Paraxial rays (close to the optical axis). Monochromatic light (no chromatic aberration). The lens is in air (or a uniform medium).


\textbf{Limitations:} Does not account for spherical aberration (rays far from the axis don't focus at the same point). Ignores chromatic aberration (different colors focus at different points). Does not apply to thick lenses, aspheric lenses, or graded-index lenses.


\section*{Mathematical Derivation}

\textbf{Step 1: Setup and Snell's Law.}

Consider a thin biconvex lens with radii of curvature $R_1$ (left surface) and $R_2$ (right surface), made of glass with refractive index $n$, embedded in air ($n_{\text{air}} \approx 1$). An object is placed at distance $d_o$ to the left of the lens. A ray from the object strikes the first surface at height $h$ from the optical axis. By Snell's law:

\begin{equation}
n_1 \sin\theta_1 = n_2 \sin\theta_2.
\end{equation}

For paraxial rays (small angles), $\sin\theta \approx \theta$, so Snell's law becomes:

\begin{equation}
n_1\,\theta_1 = n_2\,\theta_2.
\end{equation}

\textbf{Step 2: Refraction at the First Surface.}

The angle of incidence at the first surface (measured from the normal to the surface) is $\theta_1$. Using the paraxial approximation and the geometry of the curved surface, the angle of incidence is related to the object distance and height by:

\begin{equation}
\theta_1 = \frac{h}{d_o} + \frac{h}{R_1},
\end{equation}

where the first term is the angle the incident ray makes with the axis and the second comes from the surface normal tilting by $h/R_1$ (for small $h$). After refraction, the angle with respect to the axis inside the lens is:

\begin{equation}
\theta_2 = \frac{h}{R_1} - \frac{h}{d_i'},
\end{equation}

where $d_i'$ is an intermediate image distance. Applying Snell's law ($1 \cdot \theta_1 = n \cdot \theta_2$):

\begin{equation}
\frac{h}{d_o} + \frac{h}{R_1} = n\left(\frac{h}{R_1} - \frac{h}{d_i'}\right).
\end{equation}

Dividing by $h$:

\begin{equation}
\frac{1}{d_o} + \frac{1}{R_1} = \frac{n}{R_1} - \frac{n}{d_i'}. \tag{1}
\end{equation}

\textbf{Step 3: Refraction at the Second Surface.}

The ray now strikes the second surface at height approximately $h$ (thin lens: no lateral displacement). The surface has radius $R_2$ and the ray passes from glass ($n$) to air ($1$). By an analogous geometric argument:

\begin{equation}
\frac{n}{d_i'} + \frac{n}{R_2} = \frac{1}{d_i} + \frac{1}{R_2},
\end{equation}

which gives:

\begin{equation}
\frac{n}{d_i'} = \frac{1}{d_i} - \frac{n}{R_2} + \frac{1}{R_2} = \frac{1}{d_i} + \frac{1-n}{R_2}. \tag{2}
\end{equation}

\textbf{Step 4: Combining the Two Equations.}

Substitute equation (2) into equation (1):

\begin{equation}
\frac{1}{d_o} + \frac{1}{R_1} = \frac{n}{R_1} - \frac{1}{d_i} - \frac{1-n}{R_2}.
\end{equation}

Rearranging:

\begin{equation}
\frac{1}{d_o} + \frac{1}{d_i} = \frac{n-1}{R_1} - \frac{1-n}{R_2} = (n-1)\left(\frac{1}{R_1} - \frac{1}{R_2}\right).
\end{equation}

\textbf{Step 5: The Thin Lens Equation.}

We identify the right-hand side as $1/f$, where $f$ is the focal length (the lensmaker's equation):

\begin{equation}
\boxed{\frac{1}{f} = \frac{1}{d_o} + \frac{1}{d_i},}
\end{equation}

where

\begin{equation}
\frac{1}{f} = (n-1)\left(\frac{1}{R_1} - \frac{1}{R_2}\right).
\end{equation}

This is the fundamental thin lens equation. For a converging lens ($n > 1$, $R_1 > 0$, $R_2 > 0$ with $R_1 < R_2$), $f > 0$. The magnification is $M = -d_i/d_o$: a real image ($d_i > 0$) is inverted.

\section*{Characteristics of the Equation}

Magnification $M = -d_i/d_o$. A real image ($d_i > 0$) is inverted and can be projected. A virtual image ($d_i < 0$) is upright and cannot be projected. The lensmaker's equation: $1/f = (n-1)(1/R_1 - 1/R_2)$.
\section*{Solution Methods}

Ray tracing (draw three principal rays through the lens), algebraic solution using the thin lens equation, matrix optics (ABCD matrices for multi-element systems), numerical ray tracing for complex lens systems.
\section*{Verifications}

Measure object and image distances for a known lens and verify $1/f = 1/d_o + 1/d_i$. Autocollimation method: place a point source at the focal point and verify collimated output. Interferometry measures wavefront quality.
\section*{Applications}

Eyeglasses and contact lenses, camera lenses, microscope and telescope objectives, projectors, binoculars, laser focusing, fiber optic coupling, endoscopy.
\section*{What Richard Feynman Would Have Asked}

\subsection*{1. The Plain-English Translation}
\textit{``What is this equation really saying about how a lens works?''}

The Core Meaning: A lens takes light from every point on an object and redirects it so that all the light from each point converges to a single point on the image. The thin lens equation tells you where that image point is: if you know where the object is and the focal length of the lens, you know exactly where the image forms.

The Metaphor: Imagine a map where every road leads to a single city. The lens is like a magical road system that takes light from one point on an object and routes all of it to one point on the image. The focal length is how ``strong'' this routing system is---short $f$ means aggressive bending (thick lens), long $f$ means gentle bending (weak lens).

\subsection*{2. The Localized Mechanics}
\textit{``Look at a single ray of light passing through the lens. What happens to it?''}

He would press you on the refraction mechanism. At each surface of the lens, Snell's law bends the ray. The thin lens approximation assumes the lens is so thin that the ray doesn't shift laterally---it just changes direction. The focal point is defined as the point where rays initially parallel to the axis converge after passing through the lens.

The Smoothness Check: He would ask why we need the paraxial approximation. Without it, rays far from the axis don't converge to the same point---this is spherical aberration. The thin lens equation is the first-order (paraxial) approximation; higher-order corrections give the Seidel aberrations.

\subsection*{3. Testing the Extreme Limits}
\textit{``What happens when we push the lens to extreme conditions?''}

The Point-Source Limit: When $d_o \to \infty$ (object very far away), $d_i \to f$. The image forms at the focal point. This is why the Sun's image through a magnifying glass is a small spot at distance $f$---the Sun is effectively at infinity.

The Contact Lens Limit: When $d_o \to f$ from the outside, $d_i \to \infty$. The image is at infinity---the lens produces collimated (parallel) output. This is how a magnifying glass works: place the object at the focal point, and the image appears at infinity, giving maximum angular magnification.

The Diverging Limit: A concave lens ($f < 0$) always produces a virtual image between the object and the lens, no matter where the object is placed. The image is always upright and diminished---which is why concave lenses are used for nearsightedness correction.

\subsection*{4. Dimensionality and Geometry}
\textit{``What if we have two lenses instead of one?''}

He would point out that multi-element systems use the thin lens equation iteratively: the image of the first lens becomes the object of the second. The matrix method (ABCD matrices) handles this elegantly---multiply the matrices of each element to get the system matrix. The effective focal length of a two-lens system depends on the spacing between them, giving zoom capability.

\subsection*{5. Cross-Disciplinary Connection}
\textit{``Where else does this same focusing principle appear?''}

The same mathematics appears in: radio antennas (parabolic dishes focus electromagnetic waves), acoustic lenses (ultrasound focusing for medical imaging), gravitational lensing (massive galaxies bend light like lenses), and even in signal processing (lenses as Fourier transformers---the focal plane of a lens displays the spatial frequency spectrum of the input).
\section*{FAQ}

\textbf{Q: Why do cameras have multiple lens elements?}\\
\textbf{A:} A single thin lens has aberrations (spherical, chromatic). Multiple elements allow designers to cancel these aberrations---different glass types and curvatures correct for color and spherical errors, producing sharper images.

\textbf{Q: What is the difference between a converging and diverging lens?}\\
\textbf{A:} A converging (convex) lens has positive $f$ and brings parallel rays to a focus. A diverging (concave) lens has negative $f$ and spreads parallel rays apart. Your eye uses a converging lens; near-sighted (myopia) glasses use diverging lenses.
\section*{Important Bibliography}

\begin{enumerate}
\item Hecht, E., \textit{Optics}, 5th ed. (Pearson, 2017), Chapters 5--6.
\item Born, M. and Wolf, E., \textit{Principles of Optics}, 7th ed. (Cambridge University Press, 1999), Chapter 4.
\item Klein, M.V. and Furtak, T.E., \textit{Optics}, 2nd ed. (Wiley, 1986), Chapter 4.
\item Smith, W.J., \textit{Modern Optical Engineering}, 4th ed. (McGraw-Hill, 2008), Chapters 2--3.
\end{enumerate}


